% Source-derived chapter 10: Passage to   jacobian   fibrations
% Original source: no-enriques-surface-over-the-integers
\chapter{Passage to   jacobian   fibrations}
\label{no-enriques-surface-over-the-integers-chapter-10}
\source{no-enriques-surface-over-the-integers}

\begin{remark}[Source section]
\label{no-enriques-surface-over-the-integers-chapter-10-source}
\source{no-enriques-surface-over-the-integers}
\source{no-enriques-surface-over-the-integers}
This chapter reproduces the corresponding section of the retrieved
LaTeX source, including its definitions, statements, examples, and
proofs. It has not yet been translated into Lean.
\notready
\end{remark}

% The source section title was: Passage to   jacobian   fibrations
\mylabel{Passage to jacobian}

Throughout this section, $Y$ is an Enriques surface with constant Picard scheme $\Pic_{Y/k}$ over the field $k=\FF_2$.
According to Theorem \ref{family with constant pic}, 
there is at least one  genus-one fibration $\varphi:Y\ra\PP^1$. 
We now assume that this fibration is elliptic,  and let
$\phi:J\ra\PP^1$ be the resulting jacobian fibration, where $J$ is a  geometrically rational surface.
The goal of this section is to establish the following result:

\begin{theorem}
\source{no-enriques-surface-over-the-integers}
\notready
\mylabel{jacobian constant pic}
The geometrically rational surface $J$ has constant Picard scheme.
\end{theorem}

This ensures that the relation between the elliptic fibration on the Enriques surface
and its jacobian is stronger that one might expect. This crucial observation will allow us in the next sections
to use Weierstra\ss{} equations to reduce the possibilities for $J$, and thus   also for $Y$.

Choose an an algebraic closure $k\subset k^\alg$, and write $\bY=Y\otimes k^\alg$ and $\bJ=J\otimes k^\alg$
for the   base-changes.
Throughout, $\ba:\Spec(\Omega)\ra \PP^1$ denotes a  geometric point   whose image
point $a\in \PP^1$ is closed.

\begin{proposition}
\source{no-enriques-surface-over-the-integers}
\notready
\mylabel{same symbols}
For each  such  $\ba:\Spec(\Omega)\ra \PP^1$, the geometric fibers $J_\ba$ and $Y_\ba$ have the same Kodaira symbols.
If  $Y_a$ is simple,  we actually have $J_a\simeq Y_a$.
\end{proposition}

\proof
The first assertion follows from a general result of Liu, Lorenzini and Raynaud 
(\cite{Liu; Lorenzini; Raynaud 2005}, Theorem 6.6). Now suppose that $Y_a$ is simple.
In order to check $J_a\simeq Y_a$ we may base-change to $\kappa(a)=\FF_{2^\nu}$ and assume that $a\in\PP^1$ is rational.
If the fiber contains a rational point $y\in Y_a$ in the regular locus, the resulting 
effective Cartier divisor gives an identification $J\otimes R\simeq Y\otimes R$,
where $R=\O_{\PP^1,a}^h$ is the henselization, and in particular $J_a\simeq Y_a$.
It remains  to verify the existence of such a  rational point. 
For smooth fibers, this follows from Lang's result (\cite{Lang 1956}, Theorem 2).
Suppose now that $C=Y_a$ is singular, with irreducible components $C_1,\ldots, C_r$.
Each $C_i$ is birational to $\PP^1$, according to Proposition \ref{constant num over finite field}.
In case $r=1$  the description before Proposition \ref{canonical type rational points} 
reveals that the desired rational point exists.
If $r\geq 2$, we actually have $C_i\simeq \PP^1$, and
the canonical map $\Gamma(\bC)\ra\Gamma(C)$ is a graph isomorphism respecting edge labels, 
again by Proposition \ref{constant num over finite field}.
Choose a component $C_i$ corresponding to a terminal vertex of $\Gamma(C)$. Then there is exactly one
component $C_j$ intersecting $C_i$, which has $(C_i\cdot C_j)\leq 2$.
In turn, from the three rational points in $C_i=\PP^1$ at least one is contained in
the regular locus of $C_\red$. 
\qed

\medskip
We now state   several facts about the fibers of $\phi:J\ra\PP^1$, which can be formulated without referring to the Enriques surface $Y$:
\newcounter{jac}\setcounter{jac}{0}
\renewcommand{\thejac}{\rm \roman{jac}}
\newcommand{\labeljac}[1]{\refstepcounter{jac}\mylabel{#1}}
\newcommand{\refjac}[1] {{\rm (\ref{#1})}}

\begin{proposition}
\source{no-enriques-surface-over-the-integers}
\notready
\mylabel{jacobian fiber properties}
Let $\ba:\Spec(\Omega)\ra \PP^1$ be a geometric point whose image point $a\in \PP^1$ is closed.
Then the following holds:
\begin{enumerate}
\item 
\labeljac{non-rational J}
If $a\in \PP^1$ is non-rational, then the geometric fiber $J_\ba$ is irreducible.
\item
\labeljac{I0* J}
The Kodaira symbol of  $J_\ba$ is different from  $\I_0^*$.
\item
\labeljac{semistable J}
If $J_\ba$ is semistable, then  the   map $\Gamma(J_\ba)\ra\Gamma(J_a)$ is a graph isomorphism. 
\item
\labeljac{rational J}
There is at most one rational point $b\in\PP^1$ whose fiber $J_b$ is semistable or supersingular. 
\item
\labeljac{mw J}
The canonical inclusion $\MW(J/\PP^1)\subset \MW(\bJ/\bar{\PP}^1)$ 
of Mordell--Weil groups is an equality.
\end{enumerate} 
\end{proposition}

\proof
Recall from Proposition \ref{same symbols} that for all closed points $a\in \PP^1$, the fibers $J_a$ and $Y_a$ have the same Kodaira symbol.
Suppose now that $k\subset\kappa(a)$ is not an equality. Then $Y_\ba$ is   irreducible  
by Theorem \ref{geometric and schematic fiber}, which  settles \refjac{non-rational J}.
From Proposition \ref{canonical type rational points} we get \refjac{I0* J}.
Suppose next  that $J_a$   is semistable. Then $Y_a$ must be simple, by  Lemma \ref{multiple for p=2}, hence
$J_a\simeq Y_a$, again from  Proposition \ref{same symbols}. Assertion \refjac{semistable J} 
thus follows from Theorem \ref{geometric and schematic fiber}.

By Theorem \ref{family with constant pic}  there are two rational points $a\in \PP^1$ where the fiber $Y_a$ is multiple.
The corresponding $\varphi^{-1}_\ind(a)$ are neither semistable nor supersingular, again by Lemma \ref{multiple for p=2}.
So from the three rational points, there is only one $b\in\PP^1$  where $\varphi^{-1}_\ind(b)$ and hence $J_b$
can be semistable or supersingular, which gives \refjac{rational J}.

It remains to verify the last assertion \refjac{mw J}. According to the definition of Mordell--Weil groups we have
$$
\MW(J/\PP^1)=J(F)=\Pic^0_{Y_F/F}(F)=\Pic^0(Y_F),
$$
where $F$ is the function field of the projective line. The latter equality holds because $\Br(F)=0$, by Tsen's Theorem.
Choose a two-section $R\subset Y$, and consider the  surjective homomorphism
$$
\Pic(Y)\lra \Pic^0(Y_F),\quad \shL\longmapsto \shL(-nR) | Y_F,
$$
where $2n=\deg(\shL|Y_F)$. The kernel $T(Y/\PP^1)$ of this map is generated by the irreducible components 
of the closed fibers for $\varphi:Y\ra\PP^1$, together with $R$.  Consider the commutative diagram
$$
\begin{CD}
\source{no-enriques-surface-over-the-integers}
\notready
T(Y/\PP^1)	@>>> 	\Pic(Y)\\
@VVV		@VVV\\
T(\bY/\bar{\PP}^1)	@>>>	\Pic(\bY)
\end{CD}
$$
induced by base-change. By assumption,  the vertical map on the right is bijective.
According to Theorem \ref{geometric and schematic fiber}, the same holds for the vertical map on the left.
In turn, the induced vertical map  between cokernels is bijective, thus \refjac{mw J} holds.
\qed

\medskip
The number of irreducible components   of the geometric fiber $J_\ba$ depends only on the image point $a\in \PP^1$;
we denote this integer by $r_a\geq 1$. It coincides with the number of irreducible components in both $Y_\ba$ and $Y_a$.
If $a\in\PP^1$ is rational, we furthermore set
$$
n_a=
\begin{cases}
i		& \text{if $J_a$ is smooth and isomorphic to $E_i$;}\\
2r_a 	& \text{if $J_a$ is split semistable;}\\
2r_a+2	& \text{if $J_a$ is non-split semistable;}\\
2r_a+1	& \text{if $J_a$ is unstable.}
\end{cases}
$$ 
Recall from Proposition \ref{elliptic curves} that the $E_i$ denote the five elliptic curves over the prime field
$k=\FF_2$,  indexed by the order of the group of  rational points.
Also note  that $n_a=| Y_a(k)|$, and thus
$\sum n_a=25$ by  Corollary \ref{points on enriques}, where the sum runs over the three rational points $a\in\PP^1$. 
The numbers $n_a\geq 1$ are defined,  however,  
in terms of the irreducible components of the schematic and geometric fibers $J_a$ and $J_\ba$, and it is a priori not clear 
how they  are related to $|J_a(k)|$.

We now forget that $J$ comes from an elliptic Enriques surface,
and only demand that it satisfies the properties  observed above.
Theorem \ref{jacobian constant pic} becomes a consequence of the following more general statement:

\begin{theorem}
\source{no-enriques-surface-over-the-integers}
\notready
\mylabel{constant pic}
Let $J$ be a geometrically rational surface, endowed with 
an elliptic fibration $\phi:J\ra\PP^1$ that is relatively minimal and jacobian.
Suppose that conditions {\rm (i)--(v)} of Proposition \ref{jacobian fiber properties} and the equation $\sum n_a=25$ hold.
Then $J$ has constant Picard scheme.
\end{theorem}

The elliptic surface $J$ is geometrically rational, so the Picard scheme $\Pic_{Y/k}$ is  a local system
of free abelian groups. Their  rank is $\rho=10$, by the minimality of the fibration.
We have to show that the Galois group $G=\Gal(k^\alg/k)$ acts trivially on $\Pic_{J/k}(k^\alg)=\Pic(\bJ)=\ZZ^{\oplus 10}$.
Since there is no torsion part, it suffices to find a subgroup of finite index for which this holds.
The group $\Pic(\bJ)$ is generated  by a fiber, the vertical curves $\Theta$ disjoint from the zero-section,
and the sections. The former is obviously defined over $k$, and the latter  by condition \refjac{mw J}.
According to  condition \refjac{non-rational J}, each $\Theta\subset\bJ$ projects to  a fiber $J_a$ over some rational point  $a\in\PP^1$. 
Thus it suffices to prove the following:

\begin{proposition}
\source{no-enriques-surface-over-the-integers}
\notready
\mylabel{bijection irreducible components}
Let $\ba:\Spec(\Omega)\ra \PP^1$ be a geometric point. Assume that the image $a\in \PP^1$ is rational
and that the fiber $J_a$ is singular. Then  the canonical map  $\Gamma(J_\ba)\ra \Gamma(J_a)$ 
is a graph isomorphism.
\end{proposition}

\proof
The assertion  is trivial if $J_\ba$ has Kodaira symbol $\I_1$ or $\II$.
Condition \refjac{semistable J} ensures that it also holds for $\I_n$ with $n\geq 2$.
Now suppose that  $J_\ba$ has one of the Kodaira symbol $\III,\IV,\ldots,\II^*$.
Our task is to verify that the  Galois action of $G=\Gal(\FF_2^\alg/\FF_2)$ on the dual graph $\Gamma=\Gamma(J_\ba)$
is trivial.
Since it is a tree, it suffices that check that all terminal vertices are fixed.
This obviously holds for the terminal vertex $v_0\in\Gamma$ that corresponds to 
the irreducible component   that intersects the base-change of the zero-section $O\subset J$.
If there is at most one further terminal vertex, it must be fixed as well.
This already settles the cases where the Kodaira symbol is   $\III$, $\III^*$ or $\II^*$.

If the Kodaira symbol is $\I_m^*$, we must have $m\geq 1$ by condition \refjac{I0* J}.
Besides  $v_0\in\Gamma$, there are three other terminal vertices. One has a shorter distance
to $v_0$ than the others, and thus must be fixed, and we have to deal with the remaining two 
terminal vertices.
For  Kodaira symbols $\IV$ and $\IV^*$ there are also two remaining terminal vertices besides $v_0\in\Gamma$.
In all these cases, we shall check  that there is a section $P\subset J$ 
whose base-change intersects an irreducible component   $\Theta\subset J_\ba$
corresponding to one of the two remaining terminal vertices.

To achieve this, first note that $\Num(\bJ)$
is unimodular, because  the  minimal model of $\bJ$ is the projective plane.
Consider the orthogonal complement $W\subset\Num(\bJ)$ of the fiber $J_\ba$ and the zero-section $O$.
These two curves have Gram matrix $(\begin{smallmatrix}0&1\\1&-1\end{smallmatrix})$, and we infer that
$W$ is unimodular. It is generated by the vertical curves disjoint from the zero-section,
together with the combinations $(P-O)-(1+(P\cdot O))J_\ba$ for $P\in \MW(\bJ/\bar{\PP}^1)$.
The latter have self-intersection $-1-2(P\cdot O)-1$, hence  the lattice $W$ is even.
Furthermore, we consider the sublattice $W_\ba\subset W$  generated by the components $\Theta_1,\ldots,\Theta_r\subset J_\ba$ disjoint
from the zero-section. Note that this  is a root lattice, endowed with a canonical basis.

Suppose now that $J_\ba$ has Kodaira symbol $\IV$ or $\IV^*$, and that 
all sections $P\subset J$ have a base change that
passes through the component $\Theta_0\subset J_\ba$ corresponding to $v_0\in\Gamma$.
Then $W_\ba\subset W$ is an orthogonal direct summand, hence unimodular. On the other hand, this root lattice has type $A_2$ or $E_6$,
with discriminant $d=3$, contradiction.

It remains to treat the symbols $\I_m^*$ with $1\leq m\leq 4$.
Now $W_\ba$ is a root lattice of type $D_{m+4}$, which again is  not unimodular.
Our lattices and their dual lattices are related by a commutative diagram
\begin{equation}
\label{various lattices}
\source{no-enriques-surface-over-the-integers}
\begin{CD}
\source{no-enriques-surface-over-the-integers}
\notready
W_\ba	@>>>	W	@>>>	\Num(\bJ)\\
@VVV		@VVV		@VVV\\
W_\ba^*	@<<<	W^*	@<<<	\Num(\bJ)^*,
\end{CD}
\end{equation}
where the vertical maps are given by $x\mapsto (y\mapsto (x\cdot y))$. 
Seeking a contradiction, we now assume that each section $P\subset J$ passes through
the component  $\Theta_0\subset J_\ba$ corresponding to $v_0\in\Gamma$, or the
component $\Theta_1\subset J_\ba$ corresponding to the terminal vertex  $v_1\in\Gamma$  that has shorter distance to $v_0$.
The latter indeed happens, because otherwise $W_\ba\subset W$ would be an orthogonal direct summand, hence unimodular,
contradiction. Now fix some $P$ passing through $\Theta_1$.  Under the maps in \eqref{various lattices}, the combination $(P-O)-(1+(P\cdot O))J_\ba\in W$
maps to the dual basis vector  $\Theta_1^*\in W_\ba^*$. Since the maps  respect the 
intersection pairing, we conclude that $(\Theta_1^*\cdot\Theta_1^*)\in \QQ$ actually belongs to $2\ZZ$.
On the other hand, a direct computation with the root lattice $W_\ba$ of type $D_{m+4}$ reveals $(\Theta_1^*\cdot\Theta_1^*)=\pm 1$,
contradiction.
\qed

\medskip
My original arguments for the cases $\IV,\IV^*$ and $\I_m^*$ in the above proof relied on Lang's Classification \cite{Lang 2000}, 
which gives the possible  configuration of singular fibers for
$\bJ$, together with the  Oguiso--Shioda Table \cite{Oguiso; Shioda 1991}  for Mordell--Weil lattices 
and the Shioda's explicit formula \cite{Shioda 1990} for the height pairing $\langle P,Q\rangle\in\QQ$.
This  method was already used in \cite{Schroeer 2021a}, Section 15. The above more elegant
and conceptual argument was suggested by one referee.


\iffalse
Write  $b\neq c$ for the two other rational points, besides $a\in\PP^1$.
In what follows, we shall use   \emph{Lang's Classification} \cite{Lang 2000}, 
which gives the possible  configuration of singular fibers for
$\bJ$. Note that in the presence of a semistable fiber, the classification reveals that there is at most one unstable fiber.
It turns out that many cases do not occur in our situation.
For the remaining cases we use the  \emph{Oguiso--Shioda Table} \cite{Oguiso; Shioda 1991}  for Mordell--Weil lattices 
and the \emph{Shioda's explicit formula} \cite{Shioda 1990} for the \emph{height pairing} $\langle P,Q\rangle\in\QQ$  to produce
the desired section $P\subset J$.  For more details, we refer to \cite{Schroeer 2021a}, Section 15,
where the method was already employed. We now go through our four cases that the geometric fiber $J_\ba$ has Kodaira symbol
$\I_m^*$.
 
\emph{First suppose that $J_\ba$ has Kodaira symbol $\I_1^*$}, such that $n_a=13$ and $n_b+n_c=12$. By Lang's Classification,
the possible configurations of singular fibers are
$$
\I_1^*+1^4,\quad \I^*_1+21^2,\quad \I_1^*+31,\quad \I_1^*+2^2,\quad \I_1^*+4,\quad \I_1^*+\II,\quad\I_1^*+\III,\quad\I_1^*+\IV.
$$
In the first three cases we get the contradiction $n_b+n_c\leq 5+6=11$.
The fourth case does not occur thanks to condition \refjac{rational J}.
In case five the trivial lattice is the root lattice $D_5\oplus A_3$, and the Oguiso--Shioda Table gives
$\MW(J/\PP^1)=\ZZ/4\ZZ$, whereas the narrow Mordell--Weil group vanishes. Let $P\subset J$ be one of the two generators.
We have
$$
0=\langle P,P\rangle = 2+2(P\cdot O) -\contr_a(P)-\contr_b(P),
$$
and the possible non-zero values are $\contr_a(P)=1,5/4$ and $\contr_b(P)=3/4,1$.
This leaves two solutions. However, if $\contr_a(P)=\contr_b(P)=1$,
we infer that $Q=2P$ lies in the narrow Mordell--Weil group, contradiction.
Thus $\contr_a(P)=5/4$ and $\contr_b(P)=3/4$, and we infer that the base-change of $P$ intersects 
a remaining component $\Theta\subset J_\ba$ as desired.
In the remaining cases we may assume without restriction that $J_b$ is unstable and $J_c$ is smooth.
Using $n_b+n_c=12$, one sees that only the last case is possible, which was already treated.

\emph{Next, we suppose that $J_\ba$ has Kodaira symbol $\I_2^*$.}
Then $n_a=15$ and $n_b+n_c=10$.
The possible configuration of singular fibers
on $\bJ$ are
$$
\I_2^*+1^2,\quad \I_2^*+2,\quad \I_2^*.
$$
Without restriction, we may assume that $J_b$ is ordinary and $J_c$ is smooth or semi\-stable.
Using $10=n_b+n_c\leq 4+6$, we infer that $J_b=E_4$ and $J_c$ has Kodaira symbol $\tI_2$. 
The trivial lattice   is the root lattice $D_6\oplus A_1$, and we get
$\MW(J/\PP^1)=A_1^*\oplus\ZZ/2\ZZ$. Let $P\subset J$ be the
generator of the torsion subgroup.
Then
$$
0=\langle P,P\rangle = 2+2(P\cdot O) -\contr_a(P)-\contr_b(P),
$$
and the possible non-zero local contributions are $\contr_a(P)=1,3/2$ and $\contr_b(P)=1/2$.
The only solution is $\contr_a(P)=3/2$, $\contr_b(P)=1/2$ and $(P\cdot O)=0$,
hence the base-change of $P$ intersects a remaining component $\Theta\subset J_\ba$ as desired.

\emph{To continue, suppose that  $J_\ba$ has Kodaira symbol $\I_3^*$.}
Then $n_a=17$ and $n_b+n_c=8$. The only possible configuration of singular fibers 
on $\bJ$ are $\I_3^*+1$ and $\I_3^*$.
The trivial lattice for $\bJ$ is the root lattice $D_7$, and the Mordell--Weil lattice
becomes $\MW(J/\PP^1)=\langle 1/4\rangle$. Let $P\subset J$ be one of the two generators, such that 
$1/4=\langle P,P\rangle = 2+2(P\cdot O) -\contr_a(P)$. 
The   non-zero local contribution is $\contr_a(P)=1$ or $\contr_a(P)=7/4$.
The former is impossible, and we conclude as above.

\emph{It remains to treat the case that $J_\ba$ has Kodaira symbol $\I_4^*$.}
By Lang's classification, there is no other singular fiber,
so the trivial lattice for $\bJ$ is the root lattice $D_8$, and the Mordell--Weil group
must be $\MW(J/\PP^1)=\ZZ/2\ZZ$. The generator $P\subset J$ has
$0=\langle P,P\rangle = 2+2(P\cdot O) -\contr_a(P)$, and the possible non-zero local contribution is 
$\contr_a(P)=1$ or $\contr_a(P)=2$. The former is impossible, and we conclude as above.
\qed
\fi

%===========================================================
